\documentclass[a4paper,11pt]{article}

% Encodage et langue
\usepackage[T1]{fontenc}
\usepackage[utf8]{inputenc}
\usepackage[french]{babel}
\usepackage{lmodern}

% Mise en page et outils typographiques
\usepackage[a4paper,margin=2.5cm]{geometry}
\usepackage{enumitem}
\usepackage{microtype}

% Mathématiques et liens
\usepackage{amsmath,amssymb}
\usepackage[hidelinks]{hyperref}

\title{Simulation atomique 3D en C++}
\author{Notes de conception}
\date{}

\begin{document}

\maketitle

\section{Présentation du projet}

Le projet consiste à réaliser en C++ une simulation 3D simplifiée d'un atome.
La représentation graphique s'inspire du modèle de Bohr : les électrons sont
représentés autour d'un noyau constitué de protons et de neutrons.

Cette représentation est un modèle pédagogique. Elle ne décrit pas la position
réelle des électrons, qui est traitée en physique quantique par des orbitales
et des probabilités de présence.

\section{Données représentées}

\subsection{Atome}

Un atome est décrit par les informations suivantes :
\begin{itemize}[nosep]
    \item son nom et son symbole chimique, représentés par des chaînes de
          caractères (\texttt{std::string}) ;
    \item son numéro atomique, qui correspond au nombre de protons ;
    \item son nombre de neutrons ;
    \item son nombre d'électrons ;
    \item sa configuration électronique ;
    \item les coordonnées de son noyau, représentées par trois nombres
          réels : \(x\), \(y\) et \(z\).
\end{itemize}

Dans la simulation, les protons, les neutrons et les électrons peuvent être
représentés par des sphères. On peut, par exemple, distinguer les protons en
rouge, les neutrons en blanc ou en gris, et les électrons en bleu.

\subsection{Électron}

Chaque électron simulé possède notamment :
\begin{itemize}[nosep]
    \item une représentation graphique sous la forme d'une sphère ;
    \item une vitesse ;
    \item une trajectoire ;
    \item des coordonnées \(x\), \(y\) et \(z\).
\end{itemize}

\subsection{Couche électronique}

Une couche électronique correspond à un niveau principal \(n\). Pour la
visualisation, elle peut être représentée par un cercle gris centré sur le
noyau. Son diamètre dépend de la couche.

Les couches sont associées à des sous-couches \(s\), \(p\), \(d\) ou \(f\).
Le nombre maximal d'électrons d'une couche de numéro \(n\) est \(2n^2\), mais
les électrons ne remplissent pas nécessairement les couches dans l'ordre
croissant de \(n\) : ils occupent les sous-couches suivant leur énergie.

\subsection{Configuration électronique}

La configuration électronique est une liste ordonnée de sous-couches et de
leurs nombres d'électrons. Par exemple, la notation
\(\text{1s}^2\,\text{2s}^2\) indique que les sous-couches \(1s\) et \(2s\)
contiennent chacune deux électrons.

Une structure C++ peut associer à chaque sous-couche son numéro principal,
son type (\texttt{s}, \texttt{p}, \texttt{d} ou \texttt{f}) et son occupation.
Schématiquement, on peut représenter les premières occupations ainsi :
\[
    \bigl[(1s,2), (2s,2)\bigr].
\]
Il faut tenir compte de l'ordre de remplissage et de certaines exceptions.

\section{Règles de remplissage des sous-couches}

\subsection{Capacité des sous-couches}

Les sous-couches peuvent accueillir au maximum :
\begin{itemize}[nosep]
    \item sous-couche \(s\) : \(2\) électrons ;
    \item sous-couche \(p\) : \(6\) électrons ;
    \item sous-couche \(d\) : \(10\) électrons ;
    \item sous-couche \(f\) : \(14\) électrons.
\end{itemize}

\subsection{Principe de stabilité (règle de Klechkowski)}

Les sous-couches se remplissent généralement par ordre croissant d'énergie.
La règle de Klechkowski permet d'établir cet ordre à partir de la somme
\(n+\ell\), où \(n\) est le nombre quantique principal et \(\ell\) dépend du
type de sous-couche :
\[
    s \rightarrow \ell=0,\qquad
    p \rightarrow \ell=1,\qquad
    d \rightarrow \ell=2,\qquad
    f \rightarrow \ell=3.
\]

On remplit d'abord la sous-couche dont \(n+\ell\) est le plus petit. En cas
d'égalité, on remplit d'abord celle dont \(n\) est le plus petit. L'ordre
usuel commence ainsi :
\[
    1s \rightarrow 2s \rightarrow 2p \rightarrow 3s \rightarrow 3p
    \rightarrow 4s \rightarrow 3d \rightarrow 4p \rightarrow 5s
    \rightarrow 4d \rightarrow 5p \rightarrow 6s \rightarrow 4f
    \rightarrow \cdots
\]

Cet ordre général comporte des exceptions pour certains éléments, car les
niveaux d'énergie peuvent être très proches.

\subsection{Principe d'exclusion de Pauli}

Deux électrons d'un même atome ne peuvent pas avoir les quatre mêmes nombres
quantiques. En conséquence, une orbitale ne peut contenir que deux électrons
au maximum, et leurs spins doivent être opposés :
\[
    m_s=+\frac{1}{2}
    \qquad\text{et}\qquad
    m_s=-\frac{1}{2}.
\]

\subsection{Règle de Hund}

Dans des orbitales de même énergie, les électrons occupent d'abord autant
d'orbitales distinctes que possible, avec des spins parallèles. Ils commencent
à s'apparier seulement après cette répartition.

Par exemple, une sous-couche \(p\) possède trois orbitales. Avant
l'appariement, les électrons se répartissent donc individuellement dans ces
trois orbitales.

\end{document}
